\documentclass[11pt]{article}
\usepackage{fontspec}
\usepackage[french]{babel}
\usepackage{amsmath}
\usepackage{amssymb}
\usepackage{amsthm}
\usepackage{mathrsfs}
\usepackage{tikz-cd}
\usepackage{enumitem}
\usepackage{lmodern}
\usetikzlibrary{babel}

\newcommand{\printedpage}[1]{\marginpar{\footnotesize p.~#1}}
\newcommand{\msp}[1]{\marginpar{\footnotesize [#1]}}
\newcommand{\uHom}{\underline{\mathrm{Hom}}}
\newcommand{\bDelta}{\mathbf{\Delta}}
\newcommand{\ho}{\mathrm{ho}}

\begin{document}

\printedpage{4}
% head: A. Grothendieck : Dérivateurs, Ch. XIV, Édition des Universités de Montpellier II et Paris VII

\textbf{Cart (5) (Axiome d'exactitude)} Considérons deux carrés h-cartésiens
\[
X = \left(\begin{tikzcd}
X_1 \ar[d] & X_3 \ar[l] \ar[d] \\
X_0 & X_2 \ar[l]
\end{tikzcd}\right)
\qquad \text{et} \qquad
Y = \left(\begin{tikzcd}
Y_1 \ar[d] & Y_3 \ar[l] \ar[d] \\
Y_0 & Y_2 \ar[l]
\end{tikzcd}\right)
\]
et un homomorphisme $p = (p_0, p_1, p_2)$ entre les diagrammes de type
\[
\Phi \ = \
\begin{tikzcd}
1 \ar[d] & \\
0 & 2 \ar[l]
\end{tikzcd}
\]
associés (\textit{cf.} diagramme en traits pleins ci-dessous) :

\msp{page 5}
\[
\begin{tikzcd}[row sep=small, column sep=small]
& X_1 \ar[dl] \ar[dd, "p_1" near start] & & X_3 \ar[ll] \ar[dl] \ar[dd, dotted, "p_3"] \\
X_0 \ar[dd, "p_0"'] & & X_2 \ar[ll, crossing over] & \\
& Y_1 \ar[dl] & & Y_3 \ar[ll] \ar[dl] \\
Y_0 & & Y_2 \ar[ll] \ar[from=uu, crossing over, "p_2"' near end] & .
\end{tikzcd}
\]

Alors :
\begin{itemize}[leftmargin=2.5em]
\item[a)] Il existe un homomorphisme $\bar{p} : X \longrightarrow Y$ induisant $p$, \textit{i.e.} il existe une flèche $p_3$ de façon à compléter le diagramme cubique commutatif (en respectant donc la contrainte de commutativité pour les deux carrés adjacents à l'arête $p_3$).
\item[b)] Si de plus $p_0$, $p_1$, $p_2$ sont des isomorphismes, il en est de même de $p_3$ \texttt{[voir postface des éditeurs]}.
\end{itemize}

Prenant le cas où $p_0$, $p_1$, $p_2$ sont des identités, on trouve :

\textbf{Corollaire 1.} \textit{Sous les conditions de} Cart (4)\textit{, si on a} deux \textit{carrés h-cartésiens $X$, $X'$ qui complètent}
\[
\begin{tikzcd}[row sep=small]
& X_1 \ar[dl] & \\
X_0 & & , \\
& X_2 \ar[ul] &
\end{tikzcd}
\]
\textit{il existe un homomorphisme de $X$ dans $X'$ qui induise l'identité sur les sommets d'indices} 0, 1, 2\textit{, et un tel homomorphisme est toujours un isomorphisme.}

Mais cet isomorphisme n'est pas forcément unique, en d'autres termes, un automorphisme d'un carré h-cartésien peut induire l'identité sur les sommets $X_0$, $X_1$, $X_2$, sans être l'identité.
% foot: Version 7 octobre 2010

\printedpage{5}
% head: A. Grothendieck : Dérivateurs, Ch. XIV, Édition des Universités de Montpellier II et Paris VII

\msp{page 6}
\textbf{Corollaire 2.} \textit{Considérons un carré commutatif}
\[
Y = \left(\begin{tikzcd}
Y_1 \ar[d] & Y_3 \ar[l] \ar[d] \\
Y_0 & Y_2 \ar[l]
\end{tikzcd}\right) \quad ,
\]
\textit{et un objet $T$ de $\mathcal{H}$, d'où (posant} $\mathrm{Hom}(T,X) = X(T)$\textit{) un carré commutatif dans} (Ens)
\[
\begin{tikzcd}
Y_1(T) \ar[d] & Y_3(T) \ar[l] \ar[d] \\
Y_0(T) & Y_2(T) \ar[l]
\end{tikzcd} \quad ,
\]
\textit{et par suite}
\[
Y_3(T) \longrightarrow Y_1(T) \times_{Y_0(T)} Y_2(T) \quad . \tag*{$(*)$}
\]
\textit{Si $Y$ est h-cartésien, cette application} $(*)$ \textit{est} surjective \textit{quel que soit $T$.}

On applique Cart (5) à $Y$ et au carré constant $X$ de valeur $T$
\[
X = \quad
\begin{tikzcd}
T \ar[d, "\mathrm{id}"'] & T \ar[l, "\mathrm{id}"'] \ar[d, "\mathrm{id}"] \\
T & T \ar[l, "\mathrm{id}"']
\end{tikzcd} \quad ,
\]
en notant que la donnée de $p = (p_0, p_1, p_2)$ de
\[
\begin{tikzcd}
T \ar[d] & \\
T & T \ar[l]
\end{tikzcd}
\qquad\qquad \text{dans} \qquad\qquad
\begin{tikzcd}
Y_1 \ar[d] & \\
Y_0 & Y_2 \ar[l]
\end{tikzcd}
\]
revient à la donnée d'un élément du deuxième membre de $(*)$,
\[
\begin{tikzcd}
Y_1 \ar[d] & T \ar[l, "p_1"'] \ar[dl, "p_0"] \ar[d, "p_2"] \\
Y_0 & Y_2 \ar[l]
\end{tikzcd} \quad ,
\]
et que les éléments de $Y_3(T)$ qui s'envoient (par $(*)$) dans cet élément ne sont autres que les $p_3$ dont l'axiome Cart (5) affirme l'existence.

\textbf{N.B.} On notera que l'application surjective $(*)$ n'est en général \textit{pas} bijective, \textit{i.e.} le carré h-cartésien n'est en général pas cartésien.
% foot: Version 7 octobre 2010

\printedpage{6}
% head: A. Grothendieck : Dérivateurs, Ch. XIV, Édition des Universités de Montpellier II et Paris VII

\msp{page 7}
\textbf{N.B.} Il y a une notion duale de \textit{catégorie à carrés h-cocartésiens,} ou \textit{catégorie quadrangulée à droite.} Le plus souvent, les deux structures sont présentes simultanément. Il y a lieu de dégager des propriétés qui relient les deux structures duales sur une même catégorie. Ainsi, sauf erreur, dans le cas $\mathcal{H}$ additive, les deux espèces de carrés, les h-cartésiens et les h-cocartésiens, devraient \textit{coïncider.} (Il me faudra revenir là-dessus par la suite.) Dans le cas $\mathcal{H}$ ``ponctuée'', il y a l'adjonction des foncteurs $\Sigma$, $\Omega$ liés l'un à la structure quadrangulée droite, l'autre à la structure quadrangulée gauche. Et il y a, quand il y a des \underline{Hom} internes ou externes, des propriétés éventuelles de $\uHom(X,Y)$ par rapport aux quadrangles h-cocartésiens en $X$, h-cartésiens en $Y$ -- mais c'est à peine un lien entre les deux structures \ldots

\section*{2\quad Exemples}

Un premier exemple trivial, est celui où $\mathcal{H}$ est stable par produits fibrés, et où les carrés h-cartésiens sont définis comme les carrés cartésiens. Les axiomes Cart (1) à Cart (5) sont alors des propriétés bien connues, $\pm$ tautologiques de la notion de carré cartésien. Et
\msp{page 8}
en un sens, le cas général doit se ramener à ce cas particulier banal.

Le seul type d'autre cas que je connaisse, est celui où $\mathcal{H} = \mathcal{M}W^{-1}$, où $W$ est un localiseur admettant une théorie de la fibre $W$-homotopique (\textit{cf.} chapitre XII). Dans ce cas, on l'a dit, on prend comme carrés h-cartésiens dans $\mathcal{H}$ ceux qui sont isomorphes, dans $\uHom(\bDelta_1 \times \bDelta_1, \mathcal{H})$, à un carré $W$-cartésien dans $\mathcal{M}$. Encore faut-il vérifier les axiomes Cart (1) à Cart (5). Je pense que je n'ai guère de chance d'y arriver que moyennant une

\textbf{Proposition 1.} \textit{Moyennant des hypothèses sur $(\mathcal{M}, W)$ qui vont être dégagées ci-dessous, vérifiées en tous cas si $(\mathcal{M}, W)$ est quillénisable, sur la catégorie des ensembles ordonnés finis} (\textit{cf.} XIII, p. 51)\textit{, on a ceci : soit}
\[
X = \left(\begin{tikzcd}
X_1 \ar[d] & X_3 \ar[l] \ar[d] \\
X_0 & X_2 \ar[l]
\end{tikzcd}\right)
\]
\textit{un carré dans $\mathcal{M}$, et soit} $\ho_W(X)$ \textit{son image dans} $\uHom(\bDelta_1 \times \bDelta_1, \mathcal{H})$\textit{. Pour que} $\ho_W(X)$ \textit{soit h-cartésien, il} faut (\textit{et il suffit, bien sûr}) \textit{que $X$ soit $W$-cartésien} \texttt{[voir postface des éditeurs]}\textit{.}
% foot: Version 7 octobre 2010

\printedpage{7}
% head: A. Grothendieck : Dérivateurs, Ch. XIV, Édition des Universités de Montpellier II et Paris VII

\section*{Postface des éditeurs}

La proposition ci-dessus est fausse. On rappelle que si $\mathcal{M}$ est une catégorie de modèles fermée de Quillen de classe d'équivalences faibles $W$, un carré commutatif
\[
X = \left(\begin{tikzcd}
X_1 \ar[d] & X_3 \ar[l] \ar[d] \\
X_0 & X_2 \ar[l]
\end{tikzcd}\right)
\]
de $\mathcal{M}$ est $W$-cartésien si et seulement si il existe un carré \textit{cartésien}
\[
Y = \left(\begin{tikzcd}
Y_1 \ar[d] & Y_3 \ar[l] \ar[d] \\
Y_0 & Y_2 \ar[l]
\end{tikzcd}\right)
\]
avec $Y_1 \longrightarrow Y_0$ une fibration, et $Y_0$ et $Y_2$ fibrants (cette dernière condition étant superflue si la catégorie de modèles $\mathcal{M}$ est propre à droite), et un morphisme de carrés $s : X \longrightarrow Y$, $s = (s_0, s_1, s_2, s_3)$, \textit{i.e.} un cube commutatif
\[
\begin{tikzcd}[row sep=small, column sep=small]
& X_1 \ar[dl] \ar[dd, "s_1" near start] & & X_3 \ar[ll] \ar[dl] \ar[dd, "s_3"] \\
X_0 \ar[dd, "s_0"'] & & X_2 \ar[ll, crossing over] & \\
& Y_1 \ar[dl] & & Y_3 \ar[ll] \ar[dl] \\
Y_0 & & Y_2 \ar[ll] \ar[from=uu, crossing over, "s_2"' near end] & ,
\end{tikzcd}
\]
tel que $s$ soit une équivalence faible argument par argument, \textit{i.e.} $s_0, s_1, s_2, s_3 \in W$. Il est équivalent de demander que pour \textit{tout} diagramme commutatif à flèches verticales des équivalences faibles
\[
\begin{tikzcd}
X_1 \ar[r] \ar[d] & X_0 \ar[d] & X_2 \ar[l] \ar[d] \\
Y_1 \ar[r] & Y_0 & Y_2 \ar[l]
\end{tikzcd}
\]
tel que $Y_1 \longrightarrow Y_0$ soit une fibration, et $Y_0$, $Y_2$ des objets fibrants, la flèche canonique $X_3 \longrightarrow Y_1 \times_{Y_0} Y_2$ soit une équivalence faible.

Supposons maintenant que la catégorie de modèles $\mathcal{M}$ soit ponctuée, autrement dit qu'elle admette un objet $*$ à la fois initial et final, et soit $X$ un objet fibrant de $\mathcal{M}$. Choisissons une décomposition de la flèche $* \longrightarrow X$, en une équivalence faible suivie d'une fibration
\[
* \longrightarrow C \longrightarrow X \quad ,
\]
et formons le carré cartésien
% foot: Version 7 octobre 2010
\printedpage{8}
% head: A. Grothendieck : Dérivateurs, Ch. XIV, Édition des Universités de Montpellier II et Paris VII
\[
\begin{tikzcd}
C \ar[d] & \Omega X \ar[l] \ar[d] \\
X & * \ar[l]
\end{tikzcd} \quad ,
\]
où par définition $\Omega X = C \times_X *$ est l'\textit{espace des lacets} de $X$. Ce carré est $W$-cartésien. Par ailleurs, il est isomorphe dans $\uHom(\bDelta_1 \times \bDelta_1,\ \mathcal{M}W^{-1})$ au carré trivial
\[
\begin{tikzcd}
* \ar[d] & \Omega X \ar[l] \ar[d] \\
X & * \ar[l]
\end{tikzcd} \quad .
\]

Néanmoins, ce dernier n'est \textit{pas} en général $W$-cartésien. En effet, s'il l'était, l'unique flèche $s_3 : \Omega X \longrightarrow \Omega X$ de $\mathcal{M}$ rendant commutatif le cube
\[
\begin{tikzcd}[row sep=small, column sep=small]
& * \ar[dl] \ar[dd] & & \Omega X \ar[ll] \ar[dl] \ar[dd, dotted, "s_3"] \\
X \ar[dd, "\parallel"'] & & * \ar[ll, crossing over] & \\ % Audit resolved: vertical identity mark on X; see corrections.csv, PDF leaf 10.
& C \ar[dl] & & \Omega X \ar[ll] \ar[dl] \\
X & & * \ar[ll] \ar[from=uu, crossing over] &
\end{tikzcd}
\]
serait une équivalence faible. Or, comme la flèche $s_3$ se factorise par $* \times_X * \simeq *$, cela impliquerait que $\Omega X \to *$ est une équivalence faible. Cette condition est très restrictive. Si, par exemple, $\mathcal{M}$ est la catégorie de modèles des espaces topologiques pointés, et $X$ un CW-complexe pointé, elle équivaut à la contractibilité de la composante connexe de $X$ contenant le point base. De même, si $\mathcal{M}$ est la catégorie de modèles des complexes d'une catégorie abélienne de Grothendieck et $X$ un complexe, elle signifie que $X$ est acyclique.

Des considérations analogues montrent que les carrés $h$-cartésiens de $\mathcal{M}W^{-1}$, définis dans la section 2 comme étant les carrés commutatifs de $\mathcal{M}W^{-1}$ qui sont isomorphes dans $\uHom(\bDelta_1 \times \bDelta_1,\ \mathcal{M}W^{-1})$ à un carré $W$-cartésien, ne satisfont pas à toutes les conditions de Cart (3) et Cart (5) de la section 1. En effet, supposons toujours que $\mathcal{M}$ soit ponctuée, et gardons les notations ci-dessus. Considérons le diagramme commutatif dans $\mathcal{M}W^{-1}$
\[
\begin{tikzcd}
* \ar[d] & \Omega X \ar[l] \ar[d] & \Omega X \ar[l, "*_{\Omega X}"'] \ar[d] \\
X & * \ar[l] & * \ar[l]
\end{tikzcd} \quad ,
\]
où $*_{\Omega X}$ désigne l'unique endomorphisme de $\Omega X$ se factorisant par $*$. Il résulte de ce qui précède que le carré de gauche et le carré composé sont $h$-cartésiens. En revanche, le carré de droite ne l'est pas en général. En effet, si ce carré était isomorphe dans $\uHom(\bDelta_1 \times \bDelta_1,\ \mathcal{M}W^{-1})$ à un carré $W$-cartésien, la flèche $*_{\Omega X}$ serait un isomorphisme
% foot: Version 7 octobre 2010

\end{document}
